\subsection{Repeated reflections at finite scaled momentum}
\label{sec:lovelock-scaled-momentum}

At fixed \(\nu\), let \(N\to\infty\) with
\(\kappa=k^2/(4N^{2/3})\) fixed and
\(\omega-\omega_n^0\) bounded. On this ray the coupling in
\eqref{eq:lovelock-perturbed-airy} tends to
\(g_\kappa=2^{-1/3}e^{-\ii\pi/6}\kappa\). The full connection
relation of the rotated solutions gives the normalized reflection
coefficient
\begin{equation}
 \mathfrak r(g)=\frac{S_{\rm A}(g)}{S_{\rm A}(0)}
 =-\sigma\frac{a(\sigma^{-1}g)+\sigma a(\sigma g)}{a(g)}.
 \label{eq:lovelock-scaled-reflection}
\end{equation}
The denominator is nonzero on the physical ray. If \(a(g)=0\),
the recessive solution is \(O(X)\) at zero, and integration by parts
gives
\begin{equation}
 g\int_0^\infty\frac{|F|^2}{X}\,\dd X
  =\int_0^\infty (|F'|^2+X|F|^2)\,\dd X .
\end{equation}
This would require \(g\) to be positive real.

Define the optical symbol
\(\xi=e^{2\ii\omega\mathcal Z_{\rm L}-2\pi\ii h}\), which equals one
at every \(\omega_n^0\). The ingoing and outgoing boundary waves are
proportional to
\(r^{-1}H_\nu^{(1,2)}(\sqrt{\omega^2-k^2}/r)\).
The Airy connection sets their coefficient ratio to
\(Q=\xi\,\mathfrak r(g_\kappa)\) at leading order. The additive
error from the exterior momentum term in the transport action is
\(O(k^2|w|^{-1}\varepsilon_z^{-1/2})=o(1)\), where
\(\varepsilon_z\to0\) and \(\varepsilon_z|w|^{2/3}\to\infty\).
For example, at \(\varepsilon_z=|w|^{-1/2}\) the cubic correction
to the linear turning-point potential also has vanishing phase
integral, of order \(|w|\varepsilon_z^{7/2}\).
The source and response coefficients from the small-argument
Hankel expansions then give \cite{Valach:2026,Jia:2026}
\begin{align}
 \frac{\mathcal G/\mathcal G_{\rm vac}-1}{1-e^{2\pi\ii\nu}}
 &\longrightarrow \frac{Q}{1-Q},\qquad Q=\xi\,\mathfrak r(g_\kappa),
 \notag\\
 \mathcal G_{\rm vac} &=
 \frac{\Gamma(-\nu)}{\Gamma(\nu)}e^{-\ii\pi\nu}
 \left(\frac{\omega}{2}\right)^{2\nu}
 \left(1-\frac{k^2}{\omega^2}\right)^\nu .
 \label{eq:lovelock-scaled-transfer}
\end{align}
The branches are continued from positive real timelike frequency.
This normalization uses noninteger \(\nu\) and bounded frequency
offsets away from the limiting poles. At zero momentum the gamma
reflection formula applied to \eqref{eq:detuned-g} gives the
same response with \(Q=\xi\), multiplied by a gamma envelope
tending to its vacuum value.

For the branch of the pole continued from \(\kappa=0\), the limit
of the pole shift is
\begin{equation}
 \Delta\omega_{\rm A}(\kappa)
   =\frac{\ii}{2\mathcal Z_{\rm L}}\log\mathfrak r(g_\kappa),
 \label{eq:lovelock-scaled-shift}
\end{equation}
where the log is continued from zero wherever the reflection
coefficient stays nonzero. Expanding \eqref{eq:lovelock-scaled-transfer}
at \(|Q|<1\) gives \(\mathfrak r(g_\kappa)^m\) for the coefficient of
\(\xi^m\). The square and cube therefore measure the full complex
connection amplitude at the same scaled momentum. In particular,
at \(\kappa=0.45\) the calculation gives
\(\mathfrak r=3.01667-2.21406\ii\), whose magnitude is \(3.74198\).

\begin{table}[t]
 \centering\small
 \begin{tabular}{cccccc}
 \hline
 \(\kappa\) & \(\nu\) & \(n\) &
 \(\lvert\Delta\omega_n-\Delta\omega_{\rm A}\rvert\) &
 \(\lvert c_2-1\rvert\) & \(\lvert c_3-1\rvert\) \\
 \hline
 0.2 & 0.75 & 12 & 1.544e-03 & 1.035e-02 & 1.619e-02 \\
0.2 & 0.75 & 32 & 6.246e-04 & 4.164e-03 & 6.541e-03 \\
0.2 & 0.75 & 64 & 3.220e-04 & 2.138e-03 & 3.367e-03 \\
0.2 & 1.25 & 12 & 1.446e-03 & 9.128e-03 & 1.209e-02 \\
0.2 & 1.25 & 32 & 5.649e-04 & 3.739e-03 & 5.291e-03 \\
0.2 & 1.25 & 64 & 2.908e-04 & 1.950e-03 & 2.847e-03 \\
0.45 & 0.75 & 12 & 6.102e-03 & 3.627e-02 & 4.494e-02 \\
0.45 & 0.75 & 32 & 2.469e-03 & 1.468e-02 & 1.844e-02 \\
0.45 & 0.75 & 64 & 1.272e-03 & 7.553e-03 & 9.541e-03 \\
0.45 & 1.25 & 12 & 5.535e-03 & 3.335e-02 & 3.866e-02 \\
0.45 & 1.25 & 32 & 2.293e-03 & 1.384e-02 & 1.659e-02 \\
0.45 & 1.25 & 64 & 1.198e-03 & 7.210e-03 & 8.786e-03 \\
\hline

 \end{tabular}
 \caption{Finite scaled-momentum tests of the pole shift and reflection
 harmonics. The coefficients \(c_m\) are the projected harmonics of the
 full radial response divided by \(\mathfrak r(g_\kappa)^m\).}
 \label{tab:lovelock-scaled-momentum}
\end{table}

We compute \(a(g)\) by factoring
\(F=e^{-2X^{3/2}/3}X^{-1/4}H/(2\sqrt\pi)\). For
\(v=X^{-1/2}\), the large-\(X\) expansion
\(H\sim\sum_{j\geq0}d_jv^j\) has \(d_0=1\) and recurrence
\begin{equation}
 (j+1)d_{j+1}=-g\,d_j-\frac{(2j+1)(2j-3)}{16}d_{j-2},
 \qquad d_{-2}=d_{-1}=0 .
\end{equation}
From this asymptotic data we integrate inward and extract the
constant \(a(g)\). We check both adjacent Wronskians using
\(F'(X)=-g a(g)\log X+b(g)+o(1)\) at the origin and compare with
their Airy values at \(g=0\).

Table~\ref{tab:lovelock-scaled-momentum} compares these results with
independent horizon-to-boundary matching of
\eqref{eq:lovelock-k-radial}. We track the poles from \(k=0\) and
project the response along
\(\omega=\omega_n^0+(\theta-\ii\log\chi_0)/(2\mathcal Z_{\rm L})\),
for \(-\pi<\theta<\pi\) and
\(\chi_0=1/(2|\mathfrak r(g_\kappa)|)\).
At \(n=64\), changing the matching radius in \(u=r^{-2}\) from
\(0.56\) to \(0.53\) and increasing both the number of series
terms and the number of angular samples changes the pole by
less than \(10^{-8}\) and the harmonics by less than \(0.1\%\)
relative. The errors decrease with \(n\) at both values of
\(\kappa\), reaching below \(1\%\) for both harmonics at \(n=64\).
These finite-\(\kappa\) coefficients describe the
frequency-space response near the quasinormal ray. The order-\(k^2\)
polar time-domain limit can be obtained by differentiating this
connection at \(k=0\).
